\SetPractica{Alcanos, alquenos y alquinos}{Identificación de alcanos alquenos y alquinos mediante sus propiedades fisicoquímicas}

\section{Introducción teórica}

Una forma de clasificar los hidrocarburos es hacerlo de acuerdo a sus enlaces. El nombre de alcano se le da a aquellos hidrocarburos que cuentan únicamente con enlaces sencillos, mientras que a los hidrocarburos que cuentan con enlaces dobles se les llama alquenos y se les conoce como alquinos aquellos que cuentan con enlaces triples en su composición. (Wade, 2012, p. 83)

Cada uno de estos tipos de hidrocarburos tienen diferentes propiedades físicas y químicas.

Los alcanos tienen como algunas propiedades características el ser no solubles en solventes polares a causa de su propio carácter no polar y tener una densidad menor a la del agua. También, tienen un punto de ebullición bajo en moléculas con baja cantidad de carbonos, por ejemplo, el hexano tiene un punto de ebullición de 69°C. Esto es causado porque las principales fuerzas intermoleculares que definen el punto de ebullición de los alcanos son las fuerzas de van der Waals, las cuales tienen una relación proporcional con el tamaño de la molécula (Wade, 2012, p. 91). Por último, los alcanos cuentan con una buena capacidad de volatilidad, por lo que se utilizan constantemente en combustibles. Esta última propiedad suele verse afectada únicamente por la ramificación de la cadena, siendo más volátil entre menos ramificaciones tenga. (Wade, 2012, p. 93)

Los alquenos también cuentan con propiedades físicas y químicas definidas. Por su parte, tienen un carácter polar ligeramente mayor que el de los alcanos, una de las causas para esto es que los enlaces pi que presentan los alquenos tienen una fuerza más débil que los sigma, por lo que son más polarizables. La densidad de los alquenos es similar a la de los alcanos, por lo que también son insolubles en agua. Los puntos de ebullición también inician con valores bajos, pero el aumento de este valor al aumentar la cantidad de carbonos en la molécula es mayor que en los alcanos. La volatilidad de los alquenos, al igual que en los alcanos, se ve afectada por las ramificaciones. (Wade, 2012, p.298-299)

Por último, los alquinos presentan también propiedades características. Primero, tiene un carácter relativamente no polar, y casi no son solubles en agua. Su densidad es similar a la de los alcanos y alquenos, al igual que su temperatura de ebullición, la cual es solo ligeramente mayor a la de los otros tipos de hidrocarburos anteriormente mencionados. Además, durante su proceso de combustión cuenta con la característica de contar con una flama más caliente que la de los alcanos o alquenos cuando se someten a combustión. (Wade, 2012, p. 390)

Respecto a la adición de reactivos, los hidrocarburos aquí presentados cuentan con productos similares en varios de los experimentos, más no en todos.

En el caso de la adición de un halógeno, tanto el ciclohexeno como el acetileno reaccionan atacando al halógeno en su forma molecular y tomando uno de sus átomos para usar el par enlazado de sus enlaces pi para unirlo a su molécula. Esto no sucede en el caso del hexano, puesto que este no cuenta con enlaces pi de los cuales deshacerse para unir al halógeno a su composición. (Wade, 2012, p. 344, 405)

Para la adición de permanganato de potasio, los alcanos no generan ninguna reacción, puesto que no tienen enlaces dobles para unir a sus propias moléculas el permanganato. A causa de esto, el permanganato mantiene su color púrpura. En el caso del ciclohexeno, este sí que reacciona, y lo hace usando su enlace pi para enlazar un par de átomos de oxígeno de la molécula de permanganato y formar un éster cíclico. El producto se torna un color marrón. (Wade, 2012, p. 361)


\section{Objetivos de aprendizaje}

\begin{itemize}
    \item Identificar alcanos, alquenos y alquinos mediante sus propiedades fisicoquímicas
\end{itemize}

\section{Materiales, equipos y reactivos}

\begin{table}[H]
{\small\setstretch{1.25}
\begin{tabular}{|m{0.3\linewidth}|m{0.3\linewidth}|m{0.3\linewidth}|} \hline

    \thead{Equipos} & \thead{Materiales} & \thead{Reactivos} \\ \hline
        
    {Termómetro}
    
    &
    
    {Vasos de precipitado \par}
    {Cápsulas de porcelana \par}
    {Probeta \par}
    {Mechero Bunsen \par}
    {Soporte para mechero \par}
    {Malla de asbesto \par}
    {Tubos de ensayo \par}
    {Gradilla \par}
    {Pipetas graduadas}
    
    & 
    
    {Agua destilada \par}
    {Hexano \par}
    {Ciclohexeno \par}
    {Acetileno \par}
    {Solución de (\ce{Br2}/\ce{CCl4}) \par}
    {Permanganato de potasio (\ce{KMnO4}) al 1\% \par}
    {Ácido sulfúrico concentrado (\ce{H2SO4})}
        
    \\ \hline
\end{tabular}}\end{table}

\section{Procedimientos}

\subsection{Observación del estado físico y solubilidad en agua}

    \begin{enumerate}
        \item En tres tubos de ensayo, coloca una pequeña cantidad de hexano, ciclohexeno y acetileno.
        \item Añade agua destilada a cada tubo y observa si las sustancias se disuelven en agua o no.
    \end{enumerate}
    
\subsection{Punto de ebullición}

    \begin{enumerate}
        \item Determina el punto de ebullición de cada compuesto utilizando una pequeña muestra y calentando suavemente con un mechero Bunsen.
    \end{enumerate}
    
\subsection{Punto de combustión}

    \begin{enumerate}
        \item Coloca una pequeña cantidad de cada compuesto en un crisol o cápsula de porcelana en la llama de un mechero bunsen.
        \begin{itemize}
            \item Observa el tipo de llama: Los alcanos tienden a arder con una llama azulada, mientras que los alquenos y alquinos pueden producir una llama más amarillenta o incluso humeante debido a una combustión incompleta.
        \end{itemize}
    \end{enumerate}
    
\subsection{Prueba de insaturación}

    \begin{enumerate}
        \item Añade unas gotas de solución de bromo en \ce{CCl4} a tres tubos de ensayo que contengan pequeñas muestras de hexano, ciclohexeno y acetileno.
        \begin{itemize}
            \item Observa el cambio de color: Si la sustancia es insaturada (alquenos o alquinos), el color del bromo (anaranjado) desaparecerá debido a la reacción de adición, mientras que en los alcanos no habrá cambio significativo en el color.
        \end{itemize}
    \end{enumerate}
    
\subsection{Prueba de Baeyer}

    \begin{enumerate}
        \item Añade unas gotas de solución de KMnO4 al 1\% a tres tubos de ensayo con hexano, ciclohexeno y acetileno.
        \begin{itemize}
            \item Observa el cambio de color: Los alquenos y alquinos decolorarán el \ce{KMnO4} (de púrpura a marrón) debido a la oxidación, mientras que los alcanos no reaccionan.
        \end{itemize}
    \end{enumerate}
    
\subsection{Reacción con ácido sulfúrico concentrado}

    \begin{enumerate}
        \item En tres tubos de ensayo diferentes, coloca una pequeña cantidad de hexano, ciclohexeno y acetileno.
        \item Añade unas gotas de \ce{H2SO4} concentrado a cada muestra
        \begin{itemize}
            \item Observa: Los alquenos y alquinos reaccionan con el ácido sulfúrico concentrado para formar compuestos adicionales, mientras que los alcanos no reaccionan.
        \end{itemize}
    \end{enumerate}